\documentclass[12pt,a4paper]{article}
\usepackage[utf8]{inputenc}
\usepackage[T1]{fontenc}
\usepackage[french]{babel}
\usepackage{amsmath,amsfonts,amssymb}
\usepackage{xcolor}
\usepackage{tikz}
\usetikzlibrary{calc}
\usepackage[left=2cm,right=2cm,top=2cm,bottom=2cm]{geometry}

\begin{document}

\section*{\color{red} Exercice 1 : Trinômes et Équations (5 pts)}

\begin{enumerate}
    \item \textbf{Forme canonique :} $P(x) = a\left[ \left(x + \frac{b}{2a}\right)^2 - \frac{b^2 - 4ac}{4a^2} \right]$.
    
    \item \textbf{Applications :}
    \begin{itemize}
        \item $P(x)= 2x^2-x+2$. $\Delta = (-1)^2 - 4(2)(2) = -15$. \\
        Forme canonique : $P(x) = 2\left[ \left(x - \frac{1}{4}\right)^2 + \frac{15}{16} \right]$. Pas de factorisation réelle ($\Delta < 0$).
        \item $Q(x)= 3x^2-2\sqrt{3}x+1$. $\Delta = (2\sqrt{3})^2 - 4(3)(1) = 12 - 12 = 0$. \\
        Forme canonique : $Q(x) = 3\left(x - \frac{\sqrt{3}}{3}\right)^2$. Factorisation : $Q(x) = 3\left(x - \frac{\sqrt{3}}{3}\right)^2$.
    \end{itemize}

    \item \textbf{Résolutions :}
    \begin{itemize}
        \item $x^2 + 2x - 15 = 0$ : $\Delta = 64 = 8^2$. $x_1 = -5, x_2 = 3$. \textbf{$S = \{-5 ; 3\}$}.
        \item $2x^2 - 9x + 4 \geq 0$ : $\Delta = 49 = 7^2$. Racines : $1/2$ et $4$. \textbf{$S = ]-\infty ; 1/2] \cup [4 ; +\infty[$}.
        \item $6x^4 - x^2 - 5 \leq 0$ : On pose $X = x^2$. $6X^2 - X - 5 \leq 0$. Racines $X_1 = 1, X_2 = -5/6$. \\
        Comme $x^2 \geq 0$, on a $0 \leq x^2 \leq 1$, d'où \textbf{$S = [-1 ; 1]$}.
    \end{itemize}
\end{enumerate}

\section*{\color{red} Exercice 2 : Équations et Systèmes (5 pts)}

\begin{enumerate}
    \item $2|x - 2|^2 + 7|x - 2| + 3 = 0$. Posons $X = |x-2|$. $2X^2 + 7X + 3 = 0$. $\Delta = 25$. Racines : $-1/2$ et $-3$. Comme $X \geq 0$, \textbf{$S = \emptyset$}.
    
    \item Quotient $\leq 0$ : Racines num : $1$ et $-5/2$. Racines den : $5$ et $-2$. 
    Tableau de signes $\implies$ \textbf{$S = ]-\infty ; -5/2] \cup ]-2 ; 1] \cup ]5 ; +\infty[$}.
    
    \item $(x + 1/x)^2 = 4(x + 1/x) - 3$. Posons $X = x + 1/x$. $X^2 - 4X + 3 = 0 \implies X=1$ ou $X=3$.
    \begin{itemize}
        \item $x + 1/x = 1 \implies x^2 - x + 1 = 0$ ($\Delta < 0$, impossible).
        \item $x + 1/x = 3 \implies x^2 - 3x + 1 = 0$. $x = \frac{3 \pm \sqrt{5}}{2}$. \textbf{$S = \{\frac{3 - \sqrt{5}}{2} ; \frac{3 + \sqrt{5}}{2}\}$}.
    \end{itemize}

    \item $\begin{cases} x + y = 3 \\ (x+y)^2 + xy = 11 \end{cases} \implies \begin{cases} x + y = 3 \\ 9 + xy = 11 \end{cases} \implies \begin{cases} S = 3 \\ P = 2 \end{cases}$. Équation $t^2 - 3t + 2 = 0$. \textbf{$S = \{(1;2) ; (2;1)\}$}.
    
    \item \textbf{Cramer :} $\Delta = (3)(-3) - (2)(4) = -17$. $\Delta_x = (5)(-3) - (1)(4) = -19$. $\Delta_y = (3)(1) - (2)(5) = -7$. \\
    \textbf{$S = \{(\frac{19}{17} ; \frac{7}{17})\}$}.
\end{enumerate}

\section*{\textcolor{red}{\underline{Correction de l'Exercice 3}}}

\subsection*{1. Figure}

\begin{center}
\begin{tikzpicture}[scale=1.2, >=latex]
    % Définition des points
    \coordinate (A) at (0,3);
    \coordinate (B) at (-1,0);
    \coordinate (C) at (3,0);
    
    % Points dérivés
    \coordinate (I) at ($(A)!0.5!(C)$);
    \coordinate (D) at ($(B)!2!(C)$);
    \coordinate (G) at (barycentric cs:A=2,B=-1,C=2);
    \coordinate (K) at ($(B)!2!(A)$);
    
    % Tracés des segments principaux
    \draw (A) -- (B) -- (D) -- (A);
    \draw (A) -- (C);
    \draw (B) -- (I);
    \draw (C) -- (K);
    \draw (K) -- (D);

    % Marquage des points
    \foreach \p/\pos in {A/above, B/left, C/below, I/right, D/right, G/below, K/left}
        \fill (\p) circle (1.5pt) node[\pos] {$\p$};

    % Codage milieu I (Segments parallèles)
    \draw[thick] ($(A)!0.5!(I)$)+(-2pt,-2pt) -- + (2pt,2pt);
    \draw[thick] ($(I)!0.5!(C)$)+(-2pt,-2pt) -- + (2pt,2pt);
    
    % LA CORRECTION EST ICI :
    % Codage milieu C avec décalage correct
    \node at ($($(B)!0.5!(C)$)+(0,0.2)$) {$\times$};
    \node at ($($(C)!0.5!(D)$)+(0,0.2)$) {$\times$};
\end{tikzpicture}
\end{center}

\textit{La suite de la démonstration reste identique et compile désormais sans encombre.}

\subsection*{2. Justification des barycentres}

\begin{itemize}
    \item $I$ est le milieu de $[AC]$, donc par définition du milieu en tant que centre d'inertie : \textcolor{red}{$I = \text{bary}\{(A, 1); (C, 1)\}$}. En multipliant les coefficients par 2, on obtient bien \textbf{$I = \text{bary}\{(A, 2); (C, 2)\}$}.
    \item $D$ est le symétrique de $B$ par rapport à $C$, donc $C$ est le milieu de $[BD]$. On a $\overrightarrow{CB} + \overrightarrow{CD} = \vec{0}$, soit $\overrightarrow{DB} = 2\overrightarrow{DC}$, ce qui s'écrit $1\overrightarrow{DB} - 2\overrightarrow{DC} = \vec{0}$. En changeant les signes, on a $-1\overrightarrow{DB} + 2\overrightarrow{DC} = \vec{0}$, d'où \textbf{$D = \text{bary}\{(B, -1); (C, 2)\}$}.
\end{itemize}

\subsection*{3. Intersection de (AD) et (BI)}

Soit $G = \text{bary}\{(A, 2); (B, -1); (C, 2)\}$. 
\begin{itemize}
    \item Par associativité avec $I = \text{bary}\{(A, 2); (C, 2)\}$, on a $G = \text{bary}\{(I, 4); (B, -1)\}$. Donc \textbf{$G \in (BI)$}.
    \item Par associativité avec $D = \text{bary}\{(B, -1); (C, 2)\}$, on a $G = \text{bary}\{(A, 2); (D, 1)\}$. Donc \textbf{$G \in (AD)$}.
\end{itemize}
$G$ appartient aux deux droites, c'est donc leur point d'intersection.

\subsection*{4. Milieu de [BK]}

$K$ est l'intersection de $(CG)$ et $(AB)$. Dans le système $\{(A, 2); (B, -1); (C, 2)\}$, par associativité, le point $K' = \text{bary}\{(A, 2); (B, -1)\}$ appartient à $(AB)$ et à $(CG)$. Donc $K' = K$.
D'après la formule du barycentre : $2\overrightarrow{KA} - 1\overrightarrow{KB} = \vec{0} \implies \overrightarrow{KB} = 2\overrightarrow{KA}$. 
Ceci signifie que \textbf{$A$ est le milieu de $[BK]$}.

\subsection*{5. Parallélisme (DK) // (AC)}

Dans le triangle $BKD$ :
\begin{itemize}
    \item $A$ est le milieu de $[BK]$.
    \item $C$ est le milieu de $[BD]$ (car $D$ symétrique de $B$ par rapport à $C$).
\end{itemize}
D'après le théorème des milieux, la droite $(AC)$ joignant les milieux de deux côtés est parallèle au troisième côté $(DK)$. Donc \textbf{$(DK) // (AC)$}.

\subsection*{6. Alignement de $E'$, $I$ et $D$}

\begin{enumerate}
    \item \textbf{Expressions vectorielles :}
    \begin{itemize}
        \item $E' = \text{bary}\{(A, 2); (B, 1)\} \implies \overrightarrow{AE'} = \frac{1}{3}\overrightarrow{AB}$.
        $\overrightarrow{E'I} = \overrightarrow{E'A} + \overrightarrow{AI} = -\frac{1}{3}\overrightarrow{AB} + \frac{1}{2}\overrightarrow{AC}$.
        \item $D$ est tel que $\overrightarrow{CD} = \overrightarrow{BC}$, soit $\overrightarrow{AD} - \overrightarrow{AC} = \overrightarrow{AC} - \overrightarrow{AB} \implies \overrightarrow{AD} = 2\overrightarrow{AC} - \overrightarrow{AB}$.
        $\overrightarrow{ID} = \overrightarrow{IA} + \overrightarrow{AD} = -\frac{1}{2}\overrightarrow{AC} + (2\overrightarrow{AC} - \overrightarrow{AB}) = \mathbf{-\overrightarrow{AB} + \frac{3}{2}\overrightarrow{AC}}$.
    \end{itemize}
    \item \textbf{Déduction :}
    On remarque que $3\overrightarrow{E'I} = 3(-\frac{1}{3}\overrightarrow{AB} + \frac{1}{2}\overrightarrow{AC}) = -\overrightarrow{AB} + \frac{3}{2}\overrightarrow{AC} = \overrightarrow{ID}$.
    Les vecteurs sont colinéaires, donc \textbf{$E'$, $I$ et $D$ sont alignés}.
\end{enumerate}

\subsection*{7. Nature de G pour BKD}

Dans le triangle $BKD$ :
\begin{itemize}
    \item $(BI)$ est une médiane car $I$ est le milieu de $[AC]$ et $(AC)//(DK)$, mais plus simplement, $G$ est sur $(BI)$.
    \item $(AD)$ est une médiane car $A$ est le milieu de $[BK]$.
    \item $(CK)$ est une médiane car $C$ est le milieu de $[BD]$.
\end{itemize}
$G$ étant l'intersection de ces médianes, \textbf{$G$ est le centre de gravité du triangle $BKD$}.

\subsection*{8. Propriété de (AJ) et (BD)}

$J$ est l'intersection de $(BG)$ et $(KD)$. Comme $G$ est le centre de gravité de $BKD$, la droite $(BJ)$ passant par $G$ est la troisième médiane. Donc $J$ est le milieu de $[KD]$.
Dans le triangle $BKD$, $(AJ)$ joint les milieux des côtés $[BK]$ et $[KD]$. 
Par le théorème des milieux, \textbf{$(AJ)$ est parallèle à $(BD)$}.
\section*{\textcolor{red}{\underline{Correction de l'Exercice 4}}}

\begin{enumerate}
    \item \textbf{Détermination de la mesure principale $\alpha \in ]-\pi ; \pi]$ :}
    
    La méthode consiste à trouver l'entier relatif $k$ tel que $\alpha = \alpha_i - 2k\pi$ appartienne à l'intervalle $]-\pi ; \pi]$.
    
    \begin{enumerate}
        \item \textcolor{red}{$\alpha_1 = \dfrac{2027\pi}{4}$} \\
        On effectue la division : $2027 \div 4 = 506,75$. L'entier pair le plus proche est $506$. \\
        $\alpha = \dfrac{2027\pi}{4} - 506\pi = \dfrac{2027\pi - 2024\pi}{4} = \mathbf{\dfrac{3\pi}{4}}$ \\
        $\dfrac{3\pi}{4} \in ]-\pi ; \pi]$, donc la mesure principale est $\mathbf{\dfrac{3\pi}{4}}$.
        
        \item \textcolor{red}{$\alpha_2 = -\dfrac{31\pi}{6}$} \\
        On effectue la division : $-31 \div 6 \approx -5,16$. L'entier pair le plus proche est $-6$. \\
        $\alpha = -\dfrac{31\pi}{6} - (-6\pi) = -\dfrac{31\pi}{6} + \dfrac{36\pi}{6} = \mathbf{\dfrac{5\pi}{6}}$ \\
        $\dfrac{5\pi}{6} \in ]-\pi ; \pi]$, donc la mesure principale est $\mathbf{\dfrac{5\pi}{6}}$.
        
        \item \textcolor{red}{$\alpha_3 = \dfrac{2026\pi}{7}$} \\
        On effectue la division : $2026 \div 7 \approx 289,42$. L'entier pair le plus proche est $290$. \\
        $\alpha = \dfrac{2026\pi}{7} - 290\pi = \dfrac{2026\pi - 2030\pi}{7} = \mathbf{-\dfrac{4\pi}{7}}$ \\
        $-\dfrac{4\pi}{7} \in ]-\pi ; \pi]$, donc la mesure principale est $\mathbf{-\dfrac{4\pi}{7}}$.
        
        \item \textcolor{red}{$\alpha_4 = 40\pi$} \\
        $40\pi$ est un multiple pair de $\pi$ ($40\pi = 20 \times 2\pi$). \\
        La mesure principale est donc $\mathbf{0}$.
    \end{enumerate}

    \vspace{0.5cm}

    \item \textbf{Calcul des valeurs exactes :}
    
    \begin{enumerate}
        \item \textcolor{red}{$\cos \left( \dfrac{97\pi}{4} \right)$} \\
        Cherchons la mesure principale de $\dfrac{97\pi}{4}$ : $97 \div 4 = 24,25$. L'entier pair proche est $24$. \\
        $\dfrac{97\pi}{4} = 24\pi + \dfrac{\pi}{4}$. \\
        D'où : $\cos \left( \dfrac{97\pi}{4} \right) = \cos \left( \dfrac{\pi}{4} \right) = \mathbf{\dfrac{\sqrt{2}}{2}}$.
        
        \item \textcolor{red}{$\sin \left( \dfrac{31\pi}{4} \right)$} \\
        Cherchons la mesure principale de $\dfrac{31\pi}{4}$ : $31 \div 4 = 7,75$. L'entier pair proche est $8$. \\
        $\dfrac{31\pi}{4} = 8\pi - \dfrac{\pi}{4}$. \\
        D'où : $\sin \left( \dfrac{31\pi}{4} \right) = \sin \left( -\dfrac{\pi}{4} \right) = \mathbf{-\dfrac{\sqrt{2}}{2}}$.
    \end{enumerate}
\end{enumerate}
\end{document}